% =====================================================================
%  CARNET D'ENTRAINEMENT — FICHE 6 — THÈME 1
%  Niveau MOYEN — Corrigé
% =====================================================================
\documentclass[11pt,a4paper]{article}
\usepackage[utf8]{inputenc}
\usepackage[T1]{fontenc}
\usepackage{lmodern}
\usepackage[french]{babel}
\usepackage{amsmath,amssymb,mathtools}
\usepackage{icomma}
\usepackage[version=4]{mhchem}
\usepackage{siunitx}
\sisetup{output-decimal-marker={,}}
\usepackage{xcolor}
\usepackage{colortbl}
\usepackage{tikz}
\usepackage[most]{tcolorbox}
\usepackage{enumitem}
\usepackage{array}
\usepackage{booktabs}
\usepackage{fancyhdr}
\usepackage[hidelinks]{hyperref}
\usetikzlibrary{arrows.meta,positioning,calc,shapes.geometric}
\usepackage[a4paper,margin=2.1cm,headheight=26pt,headsep=8pt]{geometry}

\definecolor{primary}{RGB}{146,95,160}
\definecolor{primarydark}{RGB}{92,52,104}
\definecolor{corcol}{RGB}{0,114,178}
\definecolor{astucecol}{RGB}{198,93,9}

\tcbset{boxbase/.style={enhanced,breakable,boxrule=0.9pt,arc=2.5pt,
   left=3mm,right=3mm,top=2.2mm,bottom=2.2mm,
   fonttitle=\bfseries,coltitle=white,toptitle=0.6mm,bottomtitle=0.6mm}}
\newtcolorbox[auto counter]{cor}[1][]{boxbase,colback=corcol!5,colframe=corcol,
   title={\textsc{Corrigé}~\thetcbcounter\ifx\relax#1\relax\relax\else\textmd{ —\itshape\ #1}\fi}}
\newtcolorbox{astucebox}[1][]{boxbase,colback=astucecol!8,colframe=astucecol,
   title={\textsc{Astuce}\ifx\relax#1\relax\relax\else\textmd{ : \itshape #1}\fi}}

\pagestyle{fancy}\fancyhf{}
\fancyhead[L]{\small\textcolor{primarydark}{\textbf{Fiche 6 · Thème 1} — Corrigé}}
\fancyhead[R]{\small\textcolor{primarydark}{\textsc{Moyen}}}
\fancyfoot[C]{\small\thepage}
\renewcommand{\headrulewidth}{0.8pt}
\renewcommand{\headrule}{\hbox to\headwidth{\color{primary}\leaders\hrule height \headrulewidth\hfill}}
\setlength{\parindent}{0pt}\setlength{\parskip}{3pt}
\newcommand{\rep}{\textbf{\textcolor{corcol}{Réponse : }}}

\begin{document}
\begin{center}
{\Large\bfseries\color{primarydark}Thème 1 — Corrigé du niveau Moyen}
\end{center}
\medskip

\begin{cor}[covalent ou ionique ?]
Règle : \textbf{métal + non-métal} $\Rightarrow$ ionique (fort écart d'électronégativité) ;
\textbf{deux non-métaux} $\Rightarrow$ covalent.
\begin{itemize}[leftmargin=1.6em,itemsep=1pt]
  \item \ce{KBr} : \textbf{ionique} (\ce{K} métal alcalin + \ce{Br} non-métal).
  \item \ce{CO2} : \textbf{covalente} (\ce{C} et \ce{O}, deux non-métaux).
  \item \ce{CaO} : \textbf{ionique} (\ce{Ca} métal alcalino-terreux + \ce{O} non-métal).
  \item \ce{Cl2} : \textbf{covalente} (deux atomes identiques, écart d'électronégativité nul).
  \item \ce{HCl} : \textbf{covalente} (deux non-métaux ; liaison covalente polarisée).
  \item \ce{MgF2} : \textbf{ionique} (\ce{Mg} métal + \ce{F} non-métal).
\end{itemize}
\end{cor}

\begin{cor}[repérer le donneur d'une liaison dative]
\begin{enumerate}[label=\alph*),leftmargin=2em,itemsep=1pt]
  \item Le \textbf{donneur} est l'\textbf{oxygène} de l'eau (il possède des doublets non liants) ;
        l'\textbf{accepteur} est le proton \ce{H+} (qui a une case vide, aucun électron).
  \item L'oxygène avait 2 doublets libres ; il en \emph{engage un} dans la nouvelle liaison, il
        lui en reste donc \textbf{1}.
  \item C'est une \textbf{liaison covalente de coordination (dative)}. Une fois formée, elle est
        identique aux deux autres liaisons \ce{O-H} : \ce{H3O+} est symétrique.
\end{enumerate}
\end{cor}

\begin{cor}[repérer les exceptions à l'octet]
On regarde l'atome central de chacun.
\begin{itemize}[leftmargin=1.6em,itemsep=1pt]
  \item \ce{IF5} : \textbf{n'obéit pas} — l'iode (période 5) forme 5 liaisons et garde un doublet
        libre, soit 12 électrons : \textbf{octet étendu}.
  \item \ce{BF3} : \textbf{n'obéit pas} — le bore ne s'entoure que de 6 électrons (3 liaisons) :
        \textbf{octet incomplet} (sextet).
  \item \ce{C2H4} (éthylène) : \textbf{obéit} — chaque carbone a 8 électrons (double liaison
        \ce{C=C} + 2 liaisons \ce{C-H}).
  \item \ce{SiF4} : \textbf{obéit} — le silicium forme 4 liaisons, soit 8 électrons : octet respecté.
\end{itemize}
\rep les composés qui n'obéissent pas sont \textbf{\ce{IF5}} (octet étendu) et \textbf{\ce{BF3}}
(octet incomplet). \emph{(C'est la réponse « \ce{IF5} et \ce{BF3} » de l'annale 2020.)}
\end{cor}

\begin{cor}[peut-il dépasser l'octet ?]
Seuls les éléments de la \textbf{3\textsuperscript{e} période et au-delà} (orbitales $d$
disponibles) le peuvent.
\begin{enumerate}[label=\alph*),leftmargin=2em,itemsep=1pt]
  \item \ce{N} (\ce{NF3}) : \textbf{non} — l'azote est en 2\textsuperscript{e} période, octet strict.
  \item \ce{P} (\ce{PF5}) : \textbf{oui} — le phosphore est en 3\textsuperscript{e} période
        (ici 10 électrons, 5 liaisons).
  \item \ce{S} (\ce{SF6}) : \textbf{oui} — le soufre est en 3\textsuperscript{e} période
        (ici 12 électrons, 6 liaisons).
  \item \ce{Cl} (\ce{ClO4-}) : \textbf{oui} — le chlore est en 3\textsuperscript{e} période.
\end{enumerate}
\end{cor}

\begin{cor}[compter autour de l'atome central]
On compte $2$ électrons par liaison.
\begin{enumerate}[label=\alph*),leftmargin=2em,itemsep=1pt]
  \item \ce{BF3} : $3\times 2 = \mathbf{6}$ électrons $\Rightarrow$ octet \textbf{incomplet}
        (permis pour le bore).
  \item \ce{CH4} : $4\times 2 = \mathbf{8}$ électrons $\Rightarrow$ octet \textbf{respecté}.
  \item \ce{PCl5} : $5\times 2 = \mathbf{10}$ électrons $\Rightarrow$ octet \textbf{dépassé}
        (octet étendu, \ce{P} en 3\textsuperscript{e} période).
\end{enumerate}
\end{cor}

\begin{astucebox}[deux réflexes complémentaires]
Pour trancher une exception : (1) \textbf{qui} est l'atome central (\ce{B}/\ce{Be} $\to$ peut
rester en dessous ; période $\geq 3$ $\to$ peut dépasser ; \ce{C},\ce{N},\ce{O},\ce{F} $\to$
strict) ; (2) \textbf{combien} d'électrons l'entourent réellement (liaisons$\times2$ + libres$\times2$).
\end{astucebox}

\end{document}
